\section*{Capítulo \ref{ch:g6:life-through-seasons} --- A vida ao longo das estações}

\begin{solution}{exo:g6:life-through-seasons:1}
A duração do dia (a \omterm{def:g6:exploring-our-environment:conditions}{luz}), a \omterm{def:g6:exploring-our-environment:conditions}{temperatura} e a \omterm{def:g6:exploring-our-environment:conditions}{umidade} do solo ou do ar,
que acompanha a chuva e a seca.
\end{solution}

\begin{solution}{exo:g6:life-through-seasons:2}
\omterm{def:g6:life-through-seasons:trees}{Caducifólia} --- solta todas as \omterm{def:g1:plants-around-us:parts}{folhas} no outono: carvalho, faia,
castanheiro-da-índia. \omterm{def:g6:life-through-seasons:trees}{Perenifólia} --- mantém \omterm{def:g1:plants-around-us:parts}{folhas} trabalhando o
inverno inteiro: pinheiro, azevinho.
\end{solution}

\begin{solution}{exo:g6:life-through-seasons:3}
Um ramo completo em miniatura --- \omterm{def:g1:plants-around-us:parts}{folhas} dobradas pequenininhas,
embrulhadas em escamas à prova de tempo. Ele foi construído no verão
anterior.
\end{solution}

\begin{solution}{exo:g6:life-through-seasons:4}
A anual morre por inteiro e passa o inverno como \omterm{prop:g3:flowers-fruits-seeds:transform}{sementes}. A parte de
cima do narciso morre enquanto a \omterm{def:g1:plants-around-us:plant}{planta} espera debaixo da terra como
\omterm{ex:g4:plants-without-seeds:bulbtuber}{bulbo}.
\end{solution}

\begin{solution}{exo:g6:life-through-seasons:5}
Porque os ramos dentro dos \omterm{def:g6:life-through-seasons:buds}{brotos} dele já estavam prontos desde o
verão passado; ao galhinho só faltava calor, que a água dentro de
casa fornece semanas antes do quintal.
\end{solution}

\begin{solution}{exo:g6:life-through-seasons:6}
\omterm{def:g3:animals-through-seasons:migration}{Migrar}: a andorinha. \omterm{def:g3:animals-through-seasons:hibernation}{Hibernar}: o ouriço (ou o caracol lacrado).
Continuar ativo: a raposa, a pega, o pintarroxo. Passar o inverno nas
fases jovens --- \omterm{def:g2:animals-grow-up:born}{ovo}, larva ou \omterm{ex:g2:animals-grow-up:butterfly}{crisálida}: quase todos os \omterm{prop:g3:sorting-living-things:groups}{insetos},
como o gafanhoto ou a borboleta.
\end{solution}

\begin{solution}{exo:g6:life-through-seasons:7}
Borboleta: \omterm{ex:g2:animals-grow-up:butterfly}{crisálida} debaixo de uma pedra de cumeeira. Gafanhoto:
\omterm{def:g2:animals-grow-up:born}{ovos} na terra. Joaninha: adulta dormindo amontoada no fundo da camada
de \omterm{def:g1:plants-around-us:parts}{folhas}. Caracol: lacrado na \omterm{def:g1:animals-around-us:covering}{concha}, no barranco. Andorinha: na
África, na forma migrada dela.
\end{solution}

\begin{solution}{exo:g6:life-through-seasons:8}
Para que qualquer diferença entre as quatro tabelas venha das
estações e não do levantamento --- a regra do teste justo do
\cref{met:g6:exploring-our-environment:survey}, aplicada ao longo do
tempo em vez de entre sítios.
\end{solution}

\begin{solution}{exo:g6:life-through-seasons:9}
Primeiro, muitos ``ausentes'' estão presentes em formas escondidas
--- \omterm{def:g2:animals-grow-up:born}{ovos}, \omterm{ex:g2:animals-grow-up:butterfly}{crisálidas}, adultos dormindo, \omterm{ex:g4:plants-without-seeds:bulbtuber}{bulbos} --- que um levantamento
feito só andando por aí não encontra sem cavar e levantar. Segundo,
alguns não estão perdidos, e sim em outro lugar: os \omterm{def:g3:animals-through-seasons:migration}{migradores} vão
voltar. A lista de janeiro subestima os moradores e conta os
viajantes como sumidos.
\end{solution}

\begin{solution}{exo:g6:life-through-seasons:10}
A \omterm{def:g1:plants-around-us:plant}{planta} está no estado de depósito subterrâneo dela --- viva,
esperando na forma de \omterm{ex:g4:plants-without-seeds:bulbtuber}{bulbos}. Em março os \omterm{ex:g4:plants-without-seeds:bulbtuber}{bulbos} vão mandar \omterm{def:g1:plants-around-us:parts}{folhas} e
\omterm{def:g3:flowers-fruits-seeds:parts}{flores} para cima de novo, mais cerrados se as touceiras se
multiplicaram.
\end{solution}

\begin{solution}{exo:g6:life-through-seasons:11}
Sumiram para onde, e como o quê? Os sapos passam o inverno dormentes
--- enterrados na lama do fundo da lagoa ou em abrigos úmidos em terra
---, atravessando os meses frios em torpor até a primavera chamá-los
de volta à água para se reproduzirem.
\end{solution}

\begin{solution}{exo:g6:life-through-seasons:12}
O problema do inverno, para uma \omterm{def:g1:plants-around-us:parts}{folha} que trabalha, é segurar a água:
o chão congelado quase não dá água às \omterm{def:g1:plants-around-us:parts}{raízes}, enquanto o vento e o
sol continuam puxando água para fora --- o problema de ressecamento do
tatuzinho, só que o inverno inteiro. A cera e o formato de agulha da
\omterm{def:g1:plants-around-us:parts}{folha} \omterm{def:g6:life-through-seasons:trees}{perenifólia} cortam essas perdas até um nível que a árvore
consegue bancar. Uma \omterm{def:g1:plants-around-us:parts}{folha} larga e fina não dá para vedar tão bem; a
resposta mais barata da árvore \omterm{def:g6:life-through-seasons:trees}{caducifólia} é dar as \omterm{def:g1:plants-around-us:parts}{folhas} finas dela
por perdidas todo outono e guardar a primavera nos \omterm{def:g6:life-through-seasons:buds}{brotos}.
\end{solution}

\begin{solution}{pb:g6:life-through-seasons:1}
\textbf{1.} Os \omterm{def:g6:life-through-seasons:buds}{brotos} em cada galhinho --- o depósito de inverno que
guarda os ramos da primavera seguinte, construídos no verão anterior.

\textbf{2.} Entre os quadros de março e de abril (com os dias mais
longos e as manhãs mais amenas das anotações). A árvore responde à
\omterm{def:g6:exploring-our-environment:conditions}{temperatura} que sobe e à \omterm{def:g6:exploring-our-environment:conditions}{luz} que se alonga.

\textbf{3.} \omterm{prop:g3:sorting-living-things:groups}{Insetos} polinizadores --- as abelhas acima de tudo ---
precisam visitar as \omterm{def:g3:flowers-fruits-seeds:parts}{flores}, para o \omterm{def:g3:flowers-fruits-seeds:parts}{pólen} chegar aos \omterm{def:g3:flowers-fruits-seeds:parts}{pistilos} e a
\omterm{def:g5:flower-to-fruit:steps}{fecundação} acontecer. Uma castanha é a \omterm{def:g2:from-seed-to-plant:seed}{semente} formada a partir de um
óvulo fecundado, dentro do \omterm{prop:g3:flowers-fruits-seeds:transform}{fruto} espinhoso dela.

\textbf{4.} \omterm{def:g6:life-through-seasons:trees}{Caducifólia}. A primavera seguinte dela atravessa os
quadros de inverno na forma de \omterm{def:g6:life-through-seasons:buds}{brotos} nos galhinhos nus.

\textbf{5.} Embaixo do gramado, na forma de \omterm{ex:g4:plants-without-seeds:bulbtuber}{bulbos} --- o depósito
subterrâneo da perene.

\textbf{6.} Como \omterm{prop:g3:flowers-fruits-seeds:transform}{sementes} na terra: a estratégia da anual --- as
próprias \omterm{def:g1:plants-around-us:plant}{plantas} morreram no outono.

\textbf{7.} A andorinha: na África, na forma migrada dela, seguindo
os \omterm{prop:g3:sorting-living-things:groups}{insetos} voadores. As abelhas: na colmeia, como o cacho de inverno,
quentinhas em volta da rainha.

\textbf{8.} Continuar ativa. A pega precisa achar alimento durante os
meses mais magros --- justamente o problema de que a andorinha escapa
indo embora.

\textbf{9.} Por exemplo: os \omterm{def:g6:life-through-seasons:buds}{brotos} se abrem (março--abril) quando os
dias passam de $11$ a $13$ horas e as manhãs esquentam acima de
\qty{8}{\degreeCelsius}; as \omterm{def:g1:plants-around-us:parts}{folhas} douram e caem
(outubro--novembro) quando os dias encurtam e as manhãs esfriam. Os
acontecimentos da árvore se emparelham com mudanças de \omterm{def:g6:exploring-our-environment:conditions}{condições},
quadro após quadro.

\textbf{10.} O canteiro não morreu: ele se recolheu ao depósito ---
vivo, na forma de \omterm{ex:g4:plants-without-seeds:bulbtuber}{bulbos}, embaixo do mesmo lugar. ``Sumiu'' sempre
merece: sumiu \emph{para onde}, e como \emph{o quê}?

\textbf{11.} \omterm{def:g1:plants-around-us:tree}{Galhos} nus carregando \omterm{def:g6:life-through-seasons:buds}{brotos}; provavelmente uma pega em
algum lugar; embaixo, a camada de \omterm{def:g1:plants-around-us:parts}{folhas} e o terreno nu dos narcisos
escondendo \omterm{ex:g4:plants-without-seeds:bulbtuber}{bulbos}; \omterm{prop:g3:flowers-fruits-seeds:transform}{sementes} de papoula na terra do caminho. As
apostas certas são os \omterm{def:g6:life-through-seasons:buds}{brotos} e a nudez --- eles decorrem da
maquinaria \omterm{def:g6:life-through-seasons:trees}{caducifólia} da árvore e da prova de um ano inteiro; a pega
é provável, não é garantida.

\textbf{12.} Isso mostra que o calor era a condição que faltava --- a
\omterm{def:g6:exploring-our-environment:conditions}{luz} dentro de casa em fevereiro não é mais longa que a de fora, mas a
água e o cômodo estão quentes, e isso bastou. Teste justo: dois
galhinhos da mesma árvore, os dois em água e na mesma janela, um num
cômodo aquecido e outro numa varanda fria --- uma diferença só, o
calor; se só o galhinho aquecido se enfolhar, a questão está fechada.
\end{solution}
